A partir de la revisión de literatura expuesta en el marco referencial, se plantea como hipótesis de trabajo que las características acústicas objetivas de una señal vocal —tales como la frecuencia fundamental, la relación armónico-ruido (HNR), los coeficientes mel-frequency cepstral (MFCC), el centroide espectral, la sibilancia y el rango dinámico— guardan una relación identificable con los patrones de procesamiento que los ingenieros de sonido aplican habitualmente en cadenas de mezcla vocal profesionales (ecualización correctiva, compresión, de-essing, saturación y efectos de tiempo). Si esta relación existe y puede modelarse mediante técnicas de machine learning, entonces es posible diseñar un sistema inteligente capaz de generar, a partir del análisis de una señal vocal de entrada, propuestas iniciales de cadenas de procesamiento que resulten coherentes, técnicamente válidas y percibidas como útiles por productores con distintos niveles de experiencia, sin que ello implique sustituir el criterio creativo del ingeniero de mezcla.

De esta hipótesis se desprende, además, un supuesto complementario: que dicha asistencia inteligente reducirá la incertidumbre técnica reportada por productores amateurs y semi-profesionales —documentada por Vanka et al. en [1]— en mayor medida que en usuarios profesionales, quienes ya cuentan con criterios y flujos de trabajo consolidados, y que valorarán el sistema principalmente como punto de partida ajustable más que como solución cerrada, en línea con los hallazgos de Vanka et al. [1] sobre las diferencias de expectativas entre amateurs, semi-profesionales y profesionales frente a las herramientas de IA aplicadas a la mezcla.